\Needspace{18\baselineskip}
\subsection{Finite maps between curves}
\label{subsec:finite-curve-potential-operations}

Let $r\colon C'\to C$ be a finite surjective holomorphic map of
degree $d$ between connected compact Riemann surfaces. Write
$e_x$ for its local degree at $x\in C'$, so that
$\sum_{r(x)=s}e_x=d$ for every $s\in C$.
Fix a smooth real $(1,1)$-form $\beta$ on $C$ and put
$\sigma=r^*\beta$.
Potentials in $\PSH(C',\sigma)$ need not descend to $C$.
Taking the maximum on each fiber nevertheless gives a base
potential that dominates the original potential after pullback.
Its Lelong number at a point of $C$ is the smallest Lelong
number on the fiber over that point, after division by the
local degree.

\Needspace{15\baselineskip}
\begin{lemma}\label{finite-curve:fiber-maximum}\index{Descent!finite curve fiber maximum}
For $h\in\PSH(C',\sigma)$, the function
\[
 m_h(s)=\max_{r(x)=s}h(x)
\]
belongs to $\PSH(C,\beta)$ and satisfies $h\leq m_h\circ r$.
For every $s\in C$,
\[
 \nu(m_h,s)=\min_{r(x)=s}\frac{\nu(h,x)}{e_x}.
\]
If $v\in\PSH(C,\beta)$, then $h\preceq v\circ r$ implies
$m_h\preceq v$. The analogous implication holds for the
all-weight ideal order $\preceq_{\mathcal I}$.
\end{lemma}
\begin{proof}
On a base disk choose a smooth function $a$ with $\ddc a=\beta$.
Off the branch values, $m_h+a$ is the maximum of finitely many
subharmonic functions obtained from $h+a\circ r$ by inverse
branches. It is subharmonic there and locally bounded above.
The removable-singularity theorem extends it across the branch
values, and the extension at a branch value $s$ is
$m_h(s)+a(s)$. Indeed, upper semicontinuity of $h+a\circ r$ and
properness of $r$ bound the upper limit of $m_h+a$ at $s$ by
this value. Conversely, at each point $x$ of the fiber, the
subharmonic function $h+a\circ r$ equals its upper limit over
punctured neighborhoods of $x$, which $r$ maps into punctured
neighborhoods of $s$; this gives the reverse bound. Hence
$m_h+a$ is subharmonic on the whole disk. The function $m_h$ is
not identically $-\infty$, since $h$ is finite at some point,
so $m_h\in\PSH(C,\beta)$. The inequality $h\leq m_h\circ r$
holds by definition.

Choose a coordinate $t$ centered at $s$, and coordinates $z$
on the disjoint disks over $s$, each centered at a point $x$ of
the fiber, such that $r$ is given by $t=z^{e_x}$ on the disk
centered at $x$. Then
\[
 \sup_{|t|<\delta}m_h(t)
 =\max_{r(x)=s}\sup_{|z|<\delta^{1/e_x}}h(z).
\]
Divide by $\log\delta^2<0$ and use the supremum definition
of the Lelong number in \cref{sec:analytic-types}. Each
quotient on the right tends to
$\nu(h,x)/e_x$, and the maximum becomes the minimum after
division. Smooth local corrections do not affect these limits.

An ordinary bound $h\leq v\circ r+c$ gives
$m_h\leq v+c$, by taking the maximum over each fiber. On a
curve, local modifications are isomorphisms, so the all-weight
ideal order is equivalent to the pointwise Lelong inequalities,
by \cref{thm:local-i-inputs}. The pullback formula of
\cref{curve-map:lelong-pullback}, applied to $r$ in local
coordinates, makes $h\preceq_{\mathcal I}v\circ r$ equivalent to
$\nu(h,x)\geq e_x\nu(v,r(x))$ at every $x$.
The minimum formula then gives $\nu(m_h,s)\geq\nu(v,s)$ for
every $s\in C$, that is, $m_h\preceq_{\mathcal I}v$.
\end{proof}

The pullback of measures by $r$ counts each preimage with its
local degree, so that ramification is visible at atoms.
For a continuous real function $f$ on $C'$, put
\[
 (\operatorname{Tr}_r f)(s)
       =\sum_{r(x)=s}e_xf(x).
\]
This function is continuous: in the local description $t=z^e$,
the $e$ values of $f$ over $t\ne0$ tend to its value at the
branch point. For a finite positive measure $\mu$ on $C$, the
functional $f\mapsto\int_C\operatorname{Tr}_r f\,\mathrm d\mu$ is
positive, and it is bounded because
$|\operatorname{Tr}_r f|\leq d\|f\|_\infty$. By the Riesz
representation theorem, it defines a finite positive measure
$r^*\mu$ on $C'$ with
\[
 \int_{C'}f\,\mathrm d(r^*\mu)
       =\int_C\operatorname{Tr}_r f\,\mathrm d\mu.
\]
We extend the definition to finite signed measures by linearity.

\Needspace{16\baselineskip}
\begin{lemma}\label{finite-curve:measure-pullback}\index{Polar curvature!ramified curve pullback}
For every finite signed measure $\mu$ and Borel set $B\subset C$,
\[
 \begin{gathered}
 r_*r^*\mu=d\mu,\qquad
 \|r^*\mu\|_{\mathrm{TV}}=d\|\mu\|_{\mathrm{TV}},\\
 r^*(\mathbf1_B\mu)=\mathbf1_{r^{-1}(B)}r^*\mu,
 \qquad r^*\delta_s=\sum_{r(x)=s}e_x\delta_x.
 \end{gathered}
\]
For $v\in\PSH(C,\beta)$, write
$T_v=\beta+\ddc v$, and let
$\mu_v=\mathbf1_{\{v=-\infty\}}T_v$ be the polar curvature of
$v$, as in \cref{sec:local-polar-parts}; for potentials in
$\PSH(C',\sigma)$ the background $\sigma$ replaces $\beta$. Then
\[
 T_{v\circ r}=r^*T_v,\qquad
 \mu_{v\circ r}=r^*\mu_v.
\]
Images and inverse images of polar sets under $r$ are polar.
Pullback and pushforward of positive measures commute with
taking the atomic part.
\end{lemma}
\begin{proof}
The trace of $g\circ r$ is $dg$, giving the pushforward
identity. The same weighted fiber formula for Borel test
functions follows from the monotone-class theorem and gives
the restriction and atom formulas. The positive and negative
Jordan parts of a signed measure have disjoint Borel carriers;
their pullbacks are positive with disjoint inverse-image
carriers. Each has $d$ times its original mass, which proves
the total-variation identity.

Choose a local smooth $a$ with $\ddc a=\beta$; then $v+a$
is a subharmonic potential of $T_v$. On a smaller disk $\Delta$,
its Riesz representation is the logarithmic potential
$\int_\Delta\log|t-\tau|^2\,\mathrm dT_v(\tau)$ plus a harmonic
function. In local coordinates in which $r$ is $t=z^e$, the
identity
\[
 \ddc_z\log|z^e-\tau|^2
       =\sum_{z^e=\tau}\operatorname{ord}_z(z^e-\tau)\delta_z,
\]
valid for each fixed $\tau$, shows that pulling back this
potential gives the weighted measure $r^*T_v$. Fubini's theorem
applies locally: factorization of $z^e-\tau$ into its $e$ linear
factors gives a uniform local $L^1$ bound for these logarithmic
kernels as $\tau$ varies on a smaller disk. The harmonic term
has zero Hessian after pullback. Since
$\sigma+\ddc(v\circ r)=\ddc\bigl((v+a)\circ r\bigr)$ locally, this
proves $T_{v\circ r}=r^*T_v$.
The pole set of $v\circ r$ is exactly $r^{-1}\{v=-\infty\}$,
so the restriction formula proves the polar identity.

Inverse images of polar sets are polar by composition of
local subharmonic witnesses. For images, take local witnesses
on the finitely many inverse-image disks of a small base disk.
Off the branch values, the sum of their values over the inverse
branches is subharmonic and is $-\infty$ wherever any witness
has a pole. It is locally bounded above and is not identically
$-\infty$. Its extension across the branch values is
subharmonic. Adding the finitely many branch values to the
resulting polar set still gives a polar set. This proves the
image assertion locally, hence globally.

The pullback of a nonatomic measure has no atoms, by the
restriction and atom formulas. The pushforward of a nonatomic measure has no atoms
because every fiber is finite. Atomic measures pull back and
push forward to atomic measures. Splitting a positive measure
into these two parts proves the last assertion.
\end{proof}
